% 第二卷原书第21章思考题21.1--21.7，V2扫描PDF物理页235，候选转写，exam_ready=false。
\section*{第21章\quad 麦克斯韦方程组和电磁辐射：思考题解析}
\paragraph{思考题 21.1}
真空 SI 微分方程为 \(\nabla\cdot\boldsymbol E=\rho/\epsilon_0\)、\(\nabla\cdot\boldsymbol B=0\)、\(\nabla\times\boldsymbol E=-\partial_t\boldsymbol B\)、\(\nabla\times\boldsymbol B=\mu_0\boldsymbol J+\mu_0\epsilon_0\partial_t\boldsymbol E\)。它们依次描述电荷为电场源、无磁单极子、变化磁场产生旋电场以及电流和变化电场产生旋磁场。由后两式与连续性方程得 \(c=1/\sqrt{\mu_0\epsilon_0}\)。

\paragraph{思考题 21.2}
若采用磁荷单位使 \(\nabla\cdot\boldsymbol B=\mu_0\rho_m\)，磁荷连续性为 \(\partial_t\rho_m+\nabla\cdot\boldsymbol J_m=0\)，则可写 \(\boxed{\nabla\times\boldsymbol E=-\partial_t\boldsymbol B-\mu_0\boldsymbol J_m}\)，同时 \(\nabla\cdot\boldsymbol E=\rho_e/\epsilon_0\)、\(\nabla\times\boldsymbol B=\mu_0\boldsymbol J_e+\mu_0\epsilon_0\partial_t\boldsymbol E\)。若定义另一套磁荷单位，前面的 \(\mu_0\) 因子相应改变；仅说“加磁荷”而不说明单位制是不完整的。

\paragraph{思考题 21.3}
非相对论近似下远区的辐射场为 \(\boxed{\boldsymbol E_{\mathrm{rad}}=\frac{q}{4\pi\epsilon_0 c^2R}\,[\boldsymbol n\times(\boldsymbol n\times\boldsymbol a)]_{\mathrm{ret}}}\)，\(\boldsymbol B_{\mathrm{rad}}=(1/c)\boldsymbol n\times\boldsymbol E_{\mathrm{rad}}\)。二者横向、同相、都按 \(a_\perp/R\) 衰减；静电库仑场按 \(1/R^2\) 衰减。这里忽略相对论速度和近区 \(R^{-2},R^{-3}\) 项。

\paragraph{思考题 21.4}
坡印亭矢量 \(\boxed{\boldsymbol S=\boldsymbol E\times\boldsymbol H=(\boldsymbol E\times\boldsymbol B)/\mu_0}\)，单位 \(\mathrm{W/m^2}\)。Poynting 定理 \(\partial_t u+\nabla\cdot\boldsymbol S=-\boldsymbol J\cdot\boldsymbol E\)，其中 \(u=(\epsilon_0 E^2+B^2/\mu_0)/2\)。对真空平面波 \(\boldsymbol S\) 沿传播方向；瞬时与周期平均需区分。

\paragraph{思考题 21.5}
电偶极矩 \(\boldsymbol p(t)=\boldsymbol p_0\cos\omega t\) 时，平均辐射功率为 \(\boxed{\bar P=\omega^4p_0^2/(12\pi\epsilon_0c^3)}\)，故固定偶极矩幅值随 \(\omega^4\) 增大。若振子由给定 \(q\) 与位移振幅 \(x_0\) 构成，\(p_0=qx_0\)，仍为 \(\omega^4\)；固定加速度幅度时则不必为 \(\omega^4\)。

\paragraph{思考题 21.6}
同步辐射是相对论带电粒子在磁场中受洛伦兹力而弯曲、持续加速所产生的辐射。对 \(\gamma\gg1\) 的电子，其辐射在速度方向约 \(1/\gamma\) 的窄角内前向集束，因周期运动与脉冲压缩形成宽连续频谱，且具有方向相关的偏振。辐射来自弯曲轨迹的加速度，不需要粒子在磁场中做功；粒子损失的辐射能由加速器等补充。

\paragraph{思考题 21.7}
单位体积内有 \(n\) 个同向光子，每光子能量 \(E_\gamma=h\nu\)、动量 \(p_\gamma=h\nu/c\)。总能量密度 \(w=n h\nu\)，动量密度大小 \(g=n h\nu/c\)，故 \(\boxed{g=w/c}\)，动量密度矢量 \(\boldsymbol g=\boldsymbol S/c^2\)。原式 \(p=w/c\) 中 \(p\) 指体积动量密度，不能与单个光子动量混用。
